% three_theorems.tex
% Section 3 -- the three result-density theorems of the diamond
% progression.  This is the structural spine of the paper.  All STUB
% at v0.2; Phase 2 proves them.

\label{sec:three_theorems}

\placeholder{Section body -- write during Phase~2.  State and prove
the three theorems that characterise the geometry of the diamond
progression as a function of $d$.  Each theorem stands as its own
result; together they distinguish the diamond progression from
nearby constructions.}

\subsection{Theorem 1: Linear Coordination}
\label{subsec:theorem_1}

\begin{theorem}[Linear coordination]
\label{thm:coord_2d}
For every $d \geq 2$, every vertex of $\Lambda_d$ has exactly $2d$
neighbours:
\begin{equation*}
\text{coord}(d) = 2d.
\end{equation*}
\end{theorem}

\begin{proof}
\placeholder{Proof: count the d-of-(d+1) simplex vertex offsets and
their negations.  Each of the $d$ offsets contributes one positive
and one negative edge from any site; total $2d$.  Spell this out
in coordinates once the offset set is fixed.}
\end{proof}

\paragraph{Significance.}
This is the most consequential of the three theorems.  Chemical-
diamond and face-centred-cubic-style constructions have coordination
that stays bounded (typically~$4$ or~$12$) as $d$ grows; the diamond
progression's linear coordination is the structural feature that
distinguishes it from those constructions and that underwrites the
dimensional-selection corollary in
Section~\ref{sec:dimensional_selection}.

\subsection{Theorem 2: Maximum Euclidean Hop}
\label{subsec:theorem_2}

\begin{theorem}[Maximum hop length]
\label{thm:max_hop}
For every $d \geq 2$, the longest single edge (in Euclidean distance,
under the standard projection of the regular $(d+1)$-simplex into
$\mathbb{R}^d$) has length $\sqrt{d}$.
\end{theorem}

\begin{proof}
\placeholder{Proof: the longest edge corresponds to the simplex
vertex furthest from a chosen base vertex.  Compute the Euclidean
distance in the standard projection; it works out to $\sqrt{d}$.
Spell out the projection conventions and the explicit distance
computation.}
\end{proof}

\paragraph{Significance.}
Combined with Theorem~1, this implies that as $d$ grows the lattice
becomes \emph{denser in hops per unit Euclidean distance} -- there
are $2d$ edges per vertex but only one of length $\sqrt{d}$; the rest
are shorter.  Propagation through the $d = 3$ lattice is therefore
slower (more hops per unit physical distance) than through higher-$d$
instances.  This observation is the seed of the \emph{essay}
companion ``the progression is fractal'' tracked in `raw_thoughts.md`.

\subsection{Theorem 3: Unit-Cell Volume}
\label{subsec:theorem_3}

\begin{theorem}[Unit-cell volume]
\label{thm:unit_cell_volume}
For every $d \geq 2$, the volume of the unit cell of $\Lambda_d$ is
\begin{equation*}
\text{vol}(d) = (d+1)^{(d-1)/2}.
\end{equation*}
\end{theorem}

\begin{proof}
\placeholder{Proof: standard simplex-volume formula applied to the
$(d+1)$-vertex projection in $\mathbb{R}^d$.  $(d+1)^{(d-1)/2}$
follows from the determinant of the offset matrix.  Spell out the
matrix and the determinant computation.}
\end{proof}

\paragraph{Significance.}
Unit-cell volume grows super-exponentially in $d$ -- $\text{vol}(3)
= 4 \cdot \sqrt{2} \approx 5.66$, $\text{vol}(4) = 25$,
$\text{vol}(5) = 6^2 \sqrt{6} \approx 88$.  Combined with Theorems~1
and~2, this gives the density / scaling profile that the discrete-to-
continuum limit machinery (Section~\ref{sec:continuum_limits})
must respect.

\subsection{Why These Three}
\label{subsec:why_three_theorems}

\placeholder{The three theorems characterise the diamond progression
at three distinct geometric levels: combinatorial
(Theorem~\ref{thm:coord_2d}), Euclidean
(Theorem~\ref{thm:max_hop}), and volumetric
(Theorem~\ref{thm:unit_cell_volume}).  A fourth theorem of the same
calibre would be one of: the explicit shape of the unit cell beyond
volume (boundary structure); the second-eigenvalue / spectral gap
of the adjacency operator; the diameter / girth of $\Lambda_d$
restricted to a finite torus.  Any of these can be added in a
follow-on paper without disturbing the three-theorem spine of this
one.}

\subsection{Implications for the Rest of the Paper}
\label{subsec:three_theorem_implications}

\begin{itemize}
\item Theorem~\ref{thm:coord_2d} ($\text{coord} = 2d$) is the
      key input to the dimensional-selection corollary
      (Section~\ref{sec:dimensional_selection}).
\item Theorem~\ref{thm:max_hop} (max hop $= \sqrt{d}$) sets the
      scale parameter for the rescaled lattice $\epsilon \Lambda_d$
      in the discrete-to-continuum limit
      (Section~\ref{sec:continuum_limits}).
\item Theorem~\ref{thm:unit_cell_volume} (vol $= (d+1)^{(d-1)/2}$)
      enters the function-space scaling laws
      (Section~\ref{sec:function_spaces}) and the discrete-measure
      definitions
      (Section~\ref{sec:discrete_measure_theory}).
\end{itemize}
